%!TEX root = tfg.tex
\chapter{Adaptaciones y Excepciones del Plan de Dirección}

\begin{abstract}
Este capítulo justifica las áreas de gestión que no se desarrollan como planes
independientes dentro del proyecto. La decisión se basa en el principio de proporcionalidad:
un Plan de Dirección debe aportar control real al proyecto, no introducir documentación
adicional sin impacto práctico sobre su ejecución.
\end{abstract}

\section{Criterio de proporcionalidad}
PMBOK propone un conjunto amplio de planes y áreas de conocimiento, pero también permite
adaptar sus prácticas al contexto concreto del proyecto. En \textit{Via Inferni}, el
contexto está claramente delimitado: se trata de un Trabajo de Fin de Grado desarrollado
por dos personas, sin contratación externa, sin equipo ampliado y con un número de
interesados muy reducido.

Por este motivo, no se han desarrollado planes independientes para la gestión de
interesados, recursos y adquisiciones. Estas áreas no se ignoran, sino que se tratan de
forma integrada dentro de otros mecanismos de planificación y seguimiento más adecuados a
la escala del proyecto.

\section{Gestión de interesados}
Un plan específico de gestión de interesados no resulta necesario porque el número de
partes interesadas es limitado y estable. Los interesados principales son los dos autores
del TFG, el tutor académico y el tribunal que evaluará el trabajo. No existen clientes
externos, departamentos usuarios, patrocinadores empresariales ni grupos de adopción que
requieran segmentación, análisis de influencia o estrategias diferenciadas de implicación.

La relación con estos interesados queda cubierta por otros elementos del proyecto:

\begin{itemize}
    \item la memoria documenta el avance y las decisiones relevantes;
    \item las reuniones o revisiones con el tutor permiten validar la orientación del trabajo;
    \item los objetivos, alcance y criterios de cierre definen qué se espera entregar;
    \item la comunicación interna entre los dos autores se gestiona mediante el plan de comunicaciones.
\end{itemize}

Por tanto, elaborar un plan formal de interesados duplicaría información ya gestionada por
la planificación académica y por los mecanismos de comunicación definidos.

\section{Gestión de recursos}
Tampoco se desarrolla un plan independiente de gestión de recursos, ya que el equipo del
proyecto está formado únicamente por dos personas con roles conocidos desde el inicio. No
hay incorporación progresiva de personal, gestión de equipos externos, negociación de
disponibilidad, turnos, contratación ni reasignación compleja de perfiles.

Los recursos humanos quedan controlados mediante la planificación de horas, la distribución
de funcionalidades por iteración, la revisión cruzada y el seguimiento del trabajo en
GitHub Projects. Los recursos materiales, por su parte, se contemplan dentro del plan de
costes mediante la amortización del hardware y la identificación de herramientas empleadas.

Esta integración es suficiente para el proyecto porque la complejidad de la gestión de
recursos no está en coordinar una organización amplia, sino en asegurar que el tiempo de
los dos desarrolladores se aplica de forma realista al alcance definido.

\section{Gestión de adquisiciones}
El proyecto no requiere un plan de gestión de adquisiciones porque no contempla compras,
subcontrataciones, licitaciones, proveedores externos ni contratos asociados al desarrollo.
El videojuego se construye con herramientas disponibles para el equipo, principalmente
Unity, GitHub, Discord, Clockify y LaTeX, sin que exista una dependencia económica o
contractual que deba gestionarse formalmente.

Además, los activos necesarios para el desarrollo se generan internamente o se incorporan
bajo condiciones compatibles con un proyecto académico. En caso de utilizar recursos de
terceros, su control se limita a comprobar su licencia y atribución, lo cual se considera
parte de la gestión documental y de calidad, no una adquisición en sentido PMBOK.

Por ello, un plan de adquisiciones independiente no aportaría capacidad adicional de
control. La ausencia de compras y contratos hace que esta área pueda justificarse como no
aplicable al alcance real del TFG.

\section{Resumen de la adaptación}
La omisión de estos planes no supone una carencia metodológica, sino una adaptación
razonada del marco PMBOK. El proyecto conserva los planes que ayudan a controlar alcance,
tiempo, coste, calidad, comunicaciones y riesgos, y simplifica aquellas áreas cuya gestión
formal no resulta necesaria para un equipo de dos personas desarrollando un videojuego
académico en Unity.
